\section{Задачи}

Задачи~7.1--7.4 решаются средствами абсолютной геометрии и геометрии Лобачевского, задачи~7.5--7.11~--- с помощью модели. В задачах~7.2 и~7.4 приведено по два решения: синтетическое и в модели.

\begin{task}
Выпуклый $n$-угольник $F$ разбит диагоналями, выходящими из одной вершины, на $n-2$ треугольника. Докажите, что сумма дефектов этих треугольников равна $(n-2)\pi-s(F)$, где $s(F)$~--- сумма углов многоугольника. Выведите отсюда, что эта сумма не зависит от выбора вершины.
\end{task}
\textit{Решение.} Каждый угол многоугольника либо совпадает с углом одного из треугольников, либо (в вершине, из которой проведены диагонали) складывается из углов нескольких треугольников. Других углов у треугольников нет: все их вершины~--- вершины многоугольника. Поэтому сумма углов всех треугольников равна $s(F)$, и
\[
\sum_{i=1}^{n-2}\bigl(\pi-s_i\bigr)=(n-2)\pi-\sum_{i=1}^{n-2}s_i=(n-2)\pi-s(F).
\]
Правая часть от выбора вершины не зависит. Это позволяет назвать число $\delta(F)=(n-2)\pi-s(F)$ дефектом многоугольника.

\begin{task}
На плоскости даны пучок прямых $P$ и прямая $l$, не принадлежащая пучку. Существует ли прямая пучка, перпендикулярная $l$, и сколько их, если $P$~--- пучок а)~пересекающихся, б)~параллельных, в)~расходящихся прямых?
\end{task}
\textit{Решение 1.} а) Из центра пучка на прямую $l$ можно опустить единственный перпендикуляр, он принадлежит пучку. б) Пусть $\Omega$~--- общая точка абсолюта прямых пучка; она не является концом $l$, иначе $l$ принадлежала бы пучку. Через $\Omega$ проходит единственная прямая, перпендикулярная $l$~[1]. в) Прямая пучка перпендикулярна базе $a$; если она перпендикулярна и $l$, то это общий перпендикуляр $a$ и $l$. По теореме~\ref{thm:perpendikulyar} он существует и единствен, если $a$ и $l$ расходятся, и не существует, если они пересекаются или параллельны.

\textit{Решение 2.} Пусть пучок задан точкой $V$ (предложение~\ref{prop:puchok}), а $P_l$~--- полюс прямой $l$. По лемме~\ref{lem:polyus} прямая перпендикулярна $l$ тогда и только тогда, когда её продолжение проходит через $P_l$. Значит, искомая прямая~--- это евклидова прямая $VP_l$ (точки различны, так как $l$ не принадлежит пучку), и ответ зависит только от того, пересекает ли она круг. В случае а) $V$ лежит внутри круга, и пересечение есть. В случае б) $V$ лежит на абсолюте и не является концом хорды $l$; прямая $VP_l$ не касается абсолюта в точке $V$, так как касательная в $V$ есть поляра $V$, а $P_l$ на ней не лежит (иначе $V$ лежала бы на поляре $P_l$, то есть была бы концом $l$). Значит, прямая $VP_l$ пересекает круг. В случае в) $V$~--- полюс базы $a$, прямая $VP_l$ есть прямая через полюсы $a$ и $l$, и по теореме~\ref{thm:perpendikulyar} она пересекает круг ровно тогда, когда $a$ и $l$ расходятся.

\textit{Ответ:} в случаях а) и б) такая прямая единственна, в случае в) она существует (и единственна) тогда и только тогда, когда база пучка и прямая $l$ расходятся.

\begin{task}
Найдите дефект пятиугольника $ABCDE$, если $ABCD$~--- выпуклый четырёхугольник с углами $\angle A=\angle B=\angle C=90^\circ$, $\angle D=60^\circ$, а $ADE$~--- равносторонний треугольник с углами $\alpha$, построенный во внешнюю сторону от четырёхугольника.
\end{task}
\textit{Решение.} Треугольник $ADE$ построен во внешнюю сторону от четырёхугольника, то есть лежит по другую сторону от прямой $AD$. Поэтому в вершинах $A$ и $D$ углы четырёхугольника и треугольника прилегают друг к другу по общей стороне $AD$ и складываются. Углы пятиугольника: $\angle A=\pi/2+\alpha$, $\angle B=\angle C=\pi/2$, $\angle D=\pi/3+\alpha$, $\angle E=\alpha$. Их сумма равна $11\pi/6+3\alpha$, и
\[
\delta(ABCDE)=3\pi-\Bigl(\frac{11\pi}6+3\alpha\Bigr)=\frac{7\pi}6-3\alpha.
\]
Проверка: дефект аддитивен, $\delta(ABCD)=2\pi-11\pi/6=\pi/6$, $\delta(ADE)=\pi-3\alpha$, и сумма даёт тот же ответ.

\textit{Ответ:} $\dfrac{7\pi}{6}-3\alpha$.

\begin{task}
В выпуклом четырёхугольнике $ABCD$ углы $A$ и $B$ равны. Докажите, что серединный перпендикуляр $m$ к стороне $AB$ принадлежит пучку, которому принадлежат прямые $AD$ и $BC$.
\end{task}
\textit{Решение 1.} Симметрия $\sigma$ относительно $m$ меняет местами точки $A$ и $B$. Так как $\angle DAB=\angle CBA$ и точки $C$, $D$ лежат по одну сторону от $AB$, симметрия переводит луч $AD$ в луч $BC$, а значит, меняет местами прямые $AD$ и $BC$.

Если прямые $AD$ и $BC$ пересекаются в точке $O$, то $\sigma(O)$ лежит на обеих прямых, то есть $\sigma(O)=O$, и $O\in m$. Если они параллельны, то их общая точка абсолюта переходит в себя, поэтому она является концом прямой $m$, и $m$ параллельна обеим прямым в том же направлении. Если они расходятся, то их общий перпендикуляр $PQ$ ($P\in AD$, $Q\in BC$) переходит в общий перпендикуляр тех же прямых, то есть в себя, причём $\sigma(P)=Q\ne P$. Прямая, переходящая в себя при симметрии, либо совпадает с осью симметрии, либо ей перпендикулярна. Совпасть с $m$ прямая $PQ$ не может (точки $m$ неподвижны), поэтому $PQ\perp m$, и $m$ принадлежит пучку с базой $PQ$.

\textit{Решение 2.} Пусть $V$~--- точка пересечения евклидовых прямых, содержащих хорды $AD$ и $BC$; по предложению~\ref{prop:puchok} она задаёт пучок. Симметрия $\sigma$~--- проективное преобразование, которое меняет местами эти две прямые, поэтому их точку пересечения оставляет на месте: $\sigma(V)=V$. По формуле~\eqref{eq:otrazhenie} неподвижные точки $\sigma$~--- точки прямой $m$ и её полюс $P_m$. Если $V=P_m$, то по лемме~\ref{lem:polyus} прямые $AD$ и $BC$ перпендикулярны $m$; но прямая, перпендикулярная оси, переходит при симметрии в себя, а $\sigma$ переводит $AD$ в $BC\ne AD$. Значит, $V$ лежит на прямой, содержащей хорду $m$, и $m$ принадлежит пучку $B_V$. Это одно рассуждение для всех трёх случаев.

\begin{task}
Найдите угол параллельности отрезка, который изображается в модели отрезком от центра круга до точки $(0{,}6;\,0)$.
\end{task}
\textit{Решение.} По следствию~\ref{cor:ot-centra} длина отрезка равна $x=\operatorname{arth}0{,}6=\tfrac12\ln\tfrac{1{,}6}{0{,}4}=\ln2$. По теореме~\ref{thm:lobachevsky} $\cos\Pi(x)=\thh x=0{,}6$, откуда $\Pi(x)\approx53{,}13^\circ$. Проверка по второй формуле: $\operatorname{tg}\frac{\Pi}{2}=e^{-\ln2}=\frac12$, и $2\operatorname{arctg}\frac12\approx53{,}13^\circ$.

\textit{Ответ:} $\Pi(\ln 2)=\arccos0{,}6\approx53{,}13^\circ$.

\begin{task}
Прямые изображены хордами $k_2=0{,}3$ и $k_2=-0{,}5$. Определите их взаимное расположение и найдите длину общего перпендикуляра, если он есть.
\end{task}
\textit{Решение.} Содержащие хорды евклидовы прямые параллельны, их общая точка бесконечно удалена и лежит вне круга, поэтому прямые расходятся. Обе хорды евклидово перпендикулярны вертикальному диаметру, и по лемме~\ref{lem:polyus} он перпендикулярен обеим прямым; это и есть общий перпендикуляр (теорема~\ref{thm:perpendikulyar}). Центр лежит между его концами $(0;\,0{,}3)$ и $(0;\,-0{,}5)$, поэтому по следствию~\ref{cor:ot-centra} длина равна $\operatorname{arth}0{,}3+\operatorname{arth}0{,}5\approx0{,}3095+0{,}5493=0{,}8588$.

\textit{Ответ:} прямые расходятся, длина общего перпендикуляра $\approx0{,}859$.

\begin{task}
Найдите углы и сумму углов четырёхугольника Саккери, у которого основание и боковые стороны равны $1$.
\end{task}
\textit{Решение.} По формуле~\eqref{eq:sakkeri} при $b=h=1$
\[
\cos\gamma=\frac{\sh0{,}5\cdot\sh1}{\sqrt{\ch^20{,}5+\sh^20{,}5\cdot\sh^21}}\approx\frac{0{,}5211\cdot1{,}1752}{\sqrt{1{,}2715+0{,}2715\cdot1{,}3811}}\approx0{,}4772,
\]
$\gamma\approx61{,}49^\circ$. Сумма углов равна $2\cdot90^\circ+2\gamma\approx302{,}99^\circ$, дефект $\approx57{,}01^\circ$.

\textit{Ответ:} $\angle C=\angle D\approx61{,}49^\circ$, сумма углов $\approx302{,}99^\circ$.

\begin{task}
Найдите в модели уравнение орицикла, который касается абсолюта в точке $(1,0)$ и проходит через центр круга. Найдите центр и полуоси эллипса, которым он изображается.
\end{task}
\textit{Решение.} По формуле~\eqref{eq:klein-traektoriya} с $\omega=(1,0)$ уравнение имеет вид $(1-x)^2=c^2(1-x^2-y^2)$. Подстановка $x=y=0$ даёт $c=1$, и
\[
(1-x)^2=1-x^2-y^2\iff 2x^2-2x+y^2=0\iff\frac{(x-1/2)^2}{1/4}+\frac{y^2}{1/2}=1.
\]
\textit{Ответ:} эллипс с центром $(1/2,\,0)$, полуосью $1/2$ вдоль оси $Ox$ и полуосью $1/\sqrt2$ поперёк неё. Это согласуется с п.~5.6: $r=c/\sqrt{1+c^2}=1/\sqrt2$, $r^2=1/2$, центр $(1-r^2)\omega$.

\begin{task}
При каком расстоянии от центра кубической комнаты до стены её двугранный угол равен $80^\circ$?
\end{task}
\textit{Решение.} Из формулы~\eqref{eq:kub} $t^2=\cos80^\circ/(1+\cos80^\circ)\approx0{,}1480$, $t\approx0{,}3847$, $d=\operatorname{arth}t\approx0{,}406$.

\textit{Ответ:} примерно $0{,}406$ радиуса кривизны.

\begin{task}
Комната с полом-орисферой (п.~6.4) имеет высоту $h$. При какой высоте потолок вдвое меньше пола по длине стороны? Во сколько раз он меньше при $h=1$?
\end{task}
\textit{Решение.} По теореме~\ref{thm:komnata} сторона потолка равна стороне пола, умноженной на $e^{-h}$. Условие $e^{-h}=1/2$ даёт $h=\ln2\approx0{,}693$; при $h=1$ потолок меньше в $e\approx2{,}718$ раза.

\begin{task}
Докажите, что евклидов центр эллипсоида, изображающего сферу в модели, лежит на отрезке между центром шара и изображением центра сферы и совпадает с последним только при $q=0$.
\end{task}
\textit{Решение.} По п.~5.6 евклидов центр равен $q/B$, где $B=a^2+C(1-a^2)$, $C=\ch^2\rho$. Тогда $B-1=(C-1)(1-a^2)>0$ при $\rho>0$ и $a<1$, поэтому $0<1/B<1$, и точка $q/B$ лежит на отрезке между $0$ и $q$, не совпадая с $q$ при $q\ne0$.
